% Standalone source for Overleaf. Select XeLaTeX; no local compilation was run.
% This is a mathematical specification, not a new Lean verification report.
\documentclass[UTF8,fontset=fandol,11pt]{ctexart}
\usepackage[a4paper,margin=25mm]{geometry}
\usepackage{amsmath,amssymb,amsthm,mathtools}
\usepackage{booktabs,longtable,array,tabularx}
\usepackage{xcolor}
\usepackage{enumitem}
\usepackage{hyperref}
\usepackage{xurl}
\hypersetup{colorlinks=true,linkcolor=blue!55!black,urlcolor=blue!55!black}
\setlength{\emergencystretch}{3em}
\setlist{itemsep=3pt,topsep=5pt}
\newtheorem{proposition}{命题}[section]
\newtheorem{definition}[proposition]{定义}
\newtheorem{example}[proposition]{例}
\theoremstyle{remark}
\newtheorem{remark}[proposition]{说明}
\newcommand{\N}{\mathbb N}
\newcommand{\Q}{\mathbb Q}
\newcommand{\pa}{\operatorname{par}}
\newcommand{\ch}{\operatorname{ch}}
\newcommand{\norm}{\operatorname{norm}}
\newcommand{\EF}{\operatorname{ExactFeasible}}
\newcommand{\MF}{\operatorname{MatrixFeasible}}
\newcommand{\NC}{\operatorname{HasNextCertificate}}
\newcommand{\LC}{\operatorname{HasLocalCertificate}}
\newcommand{\RV}{\operatorname{RowValid}}
\newcommand{\CC}{\operatorname{checkCertificate}}
\newcommand{\CheckNext}{\operatorname{checkCandidate}}
\newcommand{\code}[1]{\nolinkurl{#1}}
\newcommand{\src}[2]{\href{https://github.com/huam54925-cpu/uav-routing-certification/blob/5114dc8efc5b5082aefd4640b403285d129db2b2/LeanCompletion/UAV/#1.lean}{#2}}
\title{精确路由的表示、信息保留与重建契约\\
  \large Issue \#21 第一包：数学声明与假设表}
\author{项目研究说明}
\date{2026 年 11 月 6 日}
\begin{document}
\maketitle

\begin{abstract}
本文把访问副本／next 证书、显式路线族、树父边重数矩阵和带重数的模式证书接到同一模型上。
在现有公开定理的共同假设下，给出存在性等价的纸面组合；在允许有符号费用与预算的树模型中，
给出规范化单行的见证级契约 $\pi(r)=p$ 及其证明。
路线内部顺序可以在这一投影下省略，而逐路线分组不能在检查一个给定见证时被简单加总掉。
本稿的交付是数学规格与证明文本；引用的既有 Lean 定理与本文新组合的纸面命题分别标注。
本稿不宣称 Issue \#21 已完成，也不宣称已经新增、编译或审计相应 Lean 桥接模块。
\end{abstract}

\section{固定模型和记号}

设 $V$ 为有限、具有可判定相等性的顶点类型，$o\in V$ 为基地，
$d:V\to\N$ 为访问需求，$m\in\N$ 为路线数，$s:V\to\Q$ 为每次访问的服务费，
$T\in\Q$ 为每条路线的预算。基地需求 $d(o)$ 不参与覆盖条件，基地不计服务费。
费用按\emph{每次出现}计入，经过并再次访问某个内部顶点也会再次计费；
这里不采用每个不同顶点只收费一次的语义。

\src{RouteToFlow}{\code{LegalRoute}} 以非空访问序列
$[v_1,\ldots,v_k]$ 表示路线，$k\geq1$，其中没有基地，完整顶点序列为
$[o,v_1,\ldots,v_k,o]$，相邻顶点不同。因此，每条路线从基地出发一次、最后返回一次，
途中不返回基地。允许重复访问非基地顶点，也允许不同路线的完整访问序列相同。
该类型本身不包含给定图的邻接约束；图支持需由流支持条件另外保证。

记 $F_r(u,v)$ 为路线 $r$ 在有向边 $(u,v)$ 上的使用次数，
$N_r(v)$ 为 $v$ 在访问序列中的出现次数，
\[
  C_s(r)=\sum_{v\in V}s(v)N_r(v).
\]
路线族是带索引的函数 $\rho:\operatorname{Fin}(m)\to\operatorname{LegalRoute}(o)$。
对指定有向整数流 $x:V\times V\to\N$，
\src{SupportExactFeasibility}{\code{ExactFeasible}} 的含义为
\begin{equation}\label{eq:exact}
 \exists\rho:\quad
 \begin{cases}
  \sum_i N_{\rho_i}(v)=d(v),&v\neq o,\\
  C_s(\rho_i)\leq T,&\text{每个 }i,\\
  \sum_i F_{\rho_i}(u,v)=x(u,v),&\text{每对 }u,v.
 \end{cases}
\end{equation}
第三行是\emph{每条边的重数相等}，不是仅有相同支持或边使用上界。
$m=0$ 合法，量词为空时不要求凭空补一条路线。

树侧固定 \src{TreeFlowExistence}{\code{RootedTree}} $R$：其底图为有限无向简单树，
根为 $o$，父函数为 $\pa$，且 $\pa(o)=o$。
对非根父边重数 $H:V\to\N$，定义 $E_R(H)$ 为每条父子边两个方向都使用 $H(v)$ 次的流：
\[
 E_R(H)(u,v)=
 \boldsymbol 1_{v\neq o,\,\pa(v)=u}H(v)
 +\boldsymbol 1_{u\neq o,\,\pa(u)=v}H(u).
\]
根坐标 $H(o)$ 任意且被忽略；它不是一条自环的重数。
令
\begin{align}
 k_R(p)&=\sum_{w\in\ch(o)}p(w),\label{eq:rootcount}\\
 D_R(p)(v)&=
 \begin{cases}0,&v=o,\\p(v)+\sum_{w\in\ch(v)}p(w),&v\neq o,\end{cases}\label{eq:visits}\\
 c_R(p)&=\sum_{v\in V}s(v)D_R(p)(v).\label{eq:cost}
\end{align}
其中 $\ch(v)=\{w\neq o:\pa(w)=v\}$。
需求／路线数兼容条件写作
\begin{equation}\label{eq:degrees}
 \operatorname{HasDegrees}(o,d,m,E_R(H)),
 \qquad\text{等价于}\qquad
 \begin{cases}k_R(H)=m,\\D_R(H)(v)=d(v)&(v\neq o).\end{cases}
\end{equation}
这一步不能从任意矩阵或任意 $H,d,m$ 自动略去。

\section{四种表示的类型与假设}

下表中的“存在性桥接”专指目前可直接复用的公开定理，
不把接口中的符号假设解释为数学上必然不可放宽。

\small
\begin{longtable}{@{}p{.17\textwidth}p{.35\textwidth}p{.42\textwidth}@{}}
\toprule
表示 & 数据和有效性 & 支持、符号与边界条件\\
\midrule
\endfirsthead
\toprule 表示 & 数据和有效性 & 支持、符号与边界条件\\\midrule
\endhead
访问副本／next &
$\mathcal C=\{\bot\}\sqcup\bigsqcup_{v\neq o}\{v\}\times\operatorname{Fin}(d(v))$。
每个副本有后继，任务目标的任务前驱唯一；另存严格下降的自然数秩、预算势，
并要求投影流逐边等于 $x$。 &
不要求树。后继／局部配对的表示等价本身无费用符号假设。
当前公开的双向可行性桥接要求 $x(v,v)=0$、$s(v)\geq0$、$T>0$。
指定 $m$ 时必须保留 $m=\operatorname{outFlow}(x,o)$。\\[4pt]
显式路线族 &
$\rho:\operatorname{Fin}(m)\to\operatorname{LegalRoute}(o)$，满足式 \eqref{eq:exact}。
每次访问单独计数和收费。 &
可以用于一般图；精确流 $x$ 约束其支持。
定义允许任意 $s,T\in\Q$。可含重复路线，$m=0$ 允许；单条路线的任务序列非空。\\[4pt]
树父边矩阵 &
$p_i:V\to\N$，非根列和为 $H$，每行根出发次数为 $1$，正支撑向根闭合，
且 $c_R(p_i)\leq T$。 &
要求树支持。矩阵定义忽略每行根坐标；与路线的逆向存在性对应需式 \eqref{eq:degrees}。
服务费与预算均可有符号。单行见证往返要先规定 $p_i(o)=0$。\\[4pt]
模式／重数证书 &
有限表 $[(p_j,z_j)]$，$z_j\in\N$，每个存储行满足 $\RV$，
$\sum_jz_j=m$ 且非根加权列和为 $H$。 &
树模型；每行根坐标为零且 $p_j(v)\leq H(v)$。
允许重复行和零重数，但检查器仍检查零重数行。
费用、预算允许有符号。删除零重数行的压缩对任意输入仅有单向接受保持。\\
\bottomrule
\end{longtable}
\normalsize

\subsection{next 证书中哪些信息不可省}
记副本到物理顶点的映射为 $\sigma$，$\sigma(\bot)=o$。
后继 $q$ 只要求对任务目标的任务前驱唯一，多个任务可以同时返回基地；
存储的 $q(\bot)$ 不用于生成路线边。
每个没有任务前驱的任务副本增加一条基地出发边。
严格下降秩满足
\[
 \operatorname{rank}(q(c))<\operatorname{rank}(c)\quad(c\neq\bot),
\]
以排除脱离基地的任务环。预算势满足
\[
 P(\bot)=0,\qquad
 \operatorname{fee}(c)+P(q(c))\leq P(c),\qquad
 P(c)\leq T\quad\text{若 $c$ 没有任务前驱}.
\]
未加界的 \code{NextBudgetPotential} 不要求势本身非负。
预算势不能替代下降秩：零费用的任务环可以满足全部预算不等式。

\subsection{不能把无界 next 证书与有限检查候选混为一谈}
已有 \src{LocalVisitPairingFiniteChecker}{有限 next 检查器}
\code{checkCandidate}，不调用经典前驱恢复算法，直接枚举有限副本并计算投影流。
\code{checkCandidate_eq_true_iff} 对所给候选的检查与
\code{ValidFiniteCertificate} 的等价没有额外符号假设。
但有限有效性定义本身含有 $T\geq0$ 的守卫和有限秩／势网格。
因此，不能据此把它扩展为任意有符号预算下的可行性判定器。
例如，零任务、零流、$m=0,T<0$ 时，空路线族与无界 next 证书仍存在，
该有限检查器却拒绝所有候选。
当前 \code{exactFeasible_iff_count_and_exists_checkedCandidate}
仍要求无自环、$s\geq0,T>0$。

\section{共同假设下的存在性组合}

\begin{proposition}[共同范围的四表示桥接；纸面组合]\label{prop:bridge}
固定 $R,H,d,m,s,T$，令 $x=E_R(H)$。
假设式 \eqref{eq:degrees} 成立，且
\[
 \forall v\in V,\quad 0\leq s(v),\qquad 0<T.
\]
则下列命题等价：
\begin{align*}
 &\NC(o,d,x,s,T),\\
 &\EF(o,x,d,m,s,T),\\
 &\MF(R,H,m,s,T),\\
 &\exists C,\quad \CC(R,H,m,s,T,C)=\mathrm{true}.
\end{align*}
也可将第一项换成“存在被有限 next 检查器接受的候选”。
\end{proposition}

\begin{proof}
树无自环，\code{edgeFlow_supported} 给出 $x(v,v)=0$。
由 \code{HasDegrees} 在根的出度等式得到
$m=\operatorname{outFlow}(x,o)$，所以
\src{LocalVisitPairingNextCertificate}{
\code{exactFeasible_iff_count_and_hasNextCertificate}}
右侧的路线数条件已经满足，得到前两项等价。
\src{TreeBudgetMatrix}{\code{exact_iff_matrix}} 在式 \eqref{eq:degrees} 下给出第二、三项等价。
\src{TreeBudgetCertificate}{\code{matrix_iff_exists_certificate}}
给出第三、四项等价，不另加符号假设。
有限候选版本使用
\src{LocalVisitPairingFiniteChecker}{
\code{exactFeasible_iff_count_and_exists_checkedCandidate}}，
其路线数条件由同一出度等式消去。
\end{proof}

这里组合的是\emph{存在性命题}。尚未因此得到一对可执行、互逆的证书转换程序，
也没有证明任意输入证书在所有转换后具有逐字段相等的检查结果。
把该命题写入 Lean 可作为 \#21 后续的短桥接模块；本稿没有把纸面组合标为新的内核定理。

\begin{proposition}[符号假设的精确边界；纸面组合]
在相同树数据与式 \eqref{eq:degrees} 下，允许任意 $s,T\in\Q$，仍有
\[
 \EF(o,E_R(H),d,m,s,T)
 \ \Longleftrightarrow\ \MF(R,H,m,s,T)
 \ \Longleftrightarrow\ \exists C,\ \CC(R,H,m,s,T,C)=\mathrm{true}.
\]
并且这些命题中的任一项都推出 $\NC(o,d,E_R(H),s,T)$。
\end{proposition}
\begin{proof}
前两个等价直接使用上述树定理。
\src{LocalVisitPairingCompleteness}{\code{hasLocalCertificate_of_exactFeasible}}
从路线族构造局部证书时不要求费用非负或预算为正；再使用
\src{LocalVisitPairingNextCertificate}{\code{hasNextCertificate_iff_hasLocalCertificate}}
即可获得 next 证书。
\end{proof}
现有公开接口中的非负费用和正预算假设出现在\emph{从 next／局部证书重建可行路线族}
的方向。这里不把未复用到的有符号逆向定理说成不可能，也不提前声称它已经形式化。
若使用 \code{PatternCountsFeasible} 的函数式计数版本，则还需显式使用
$m=k_R(H)$：该定义省略了独立路线数参数，
\code{pattern_count_total} 从列等式推导总模式重数为 $k_R(H)$。

\section{单行的投影和重建契约}

\begin{definition}[规范化行与投影]
设 $\norm(p)(o)=0$，$\norm(p)(v)=p(v)$（$v\neq o$）。
对任意合法路线 $r$ 定义
\[
 \pi_R(r)(v)=
 \begin{cases}0,&v=o,\\F_r(\pa(v),v),&v\neq o.\end{cases}
\]
这等于对现有 \code{TreeBudgetRow.routeRow} 再应用
\code{TreeBudgetPatterns.normalize}。
对树支持的路线，因 $\pa(o)=o$ 且无自环，原 \code{routeRow} 的根坐标也已为零。
显式写规范化能避免将矩阵定义中本来任意的根坐标带入往返等式。
\end{definition}

正支撑向根闭合是
\begin{equation}\label{eq:ancestor}
 v\neq o,\quad p(v)>0,\quad\pa(v)\neq o
 \quad\Longrightarrow\quad p(\pa(v))>0.
\end{equation}
现有行谓词 $\RV(R,H,s,T,p)$ 的五项条件依次为
\begin{equation}\label{eq:rowvalid}
 p(o)=0,\qquad p(v)\leq H(v)\ (v\neq o),\qquad
 k_R(p)=1,\qquad\eqref{eq:ancestor},\qquad c_R(p)\leq T.
\end{equation}
现有可执行 \code{cost} 与 $c_R$ 相等，且有线性形式
\begin{equation}\label{eq:linear}
 c_R(p)=\sum_{v\neq o}
 \bigl(s(v)+\boldsymbol 1_{\pa(v)\neq o}s(\pa(v))\bigr)p(v).
\end{equation}
式 \eqref{eq:visits} 和 \eqref{eq:linear} 已由
\src{TreeBudgetRow}{\code{routeRow_counts}、\code{rowCost_linear}}
等定理支持，不把它们重复计为本包的新成果。

\begin{proposition}[规范化单行的见证级重建；纸面推论]\label{prop:section}
若 $p:V\to\N$ 满足 $p(o)=0$、$k_R(p)=1$ 和式 \eqref{eq:ancestor}，
则存在合法路线 $r$，满足
\begin{equation}\label{eq:section}
 F_r=E_R(p),\qquad \pi_R(r)=p,\qquad
 N_r(v)=D_R(p)(v)\ (v\in V),\qquad C_s(r)=c_R(p).
\end{equation}
此处 $s:V\to\Q$ 任意。若进一步有 $\RV(R,H,s,T,p)$，则所构造的路线满足预算 $T$。
\end{proposition}
\begin{proof}
将 $k_R(p)=1$ 通过 \code{rootCount_eq} 换成 \code{rootDegree R p = 1}。
\src{TreeBudgetRow}{\code{single_row_realization}}
给出 $r$，其逐边流恰为 $E_R(p)$，访问数恰为 $D_R(p)$。
对 $v\neq o$，\code{edgeFlow_parent} 给出
\[
 \pi_R(r)(v)=F_r(\pa(v),v)=E_R(p)(\pa(v),v)=p(v).
\]
根坐标由 $p(o)=0$ 恢复，故函数外延相等成立。
逐顶点计费并求和，由 \code{service_of_counts} 得到成本等式。
预算来自 $c_R(p)\leq T$。
\end{proof}

这说明可对每个满足条件的 $p$ \emph{选取}一个实现，记作
$\operatorname{realize}_R(p)$，并满足
\[
 \boxed{\ \pi_R\bigl(\operatorname{realize}_R(p)\bigr)=p\ }.
\]
这是投影的右逆／一个截面，而非双向逆。
现有 \code{single_row_realization} 给出存在性证明，
用选择公理选取见证是非计算性的；本文没有指定可执行遍历顺序，也没有实现命名为
\code{realize} 的 Lean 函数。
若去掉 $p(o)=0$，同一证明只能得到 $\pi_R(r)=\norm(p)$。
列上界 $p\leq H$ 和预算不是结构性重建的必要前提；它们用于接入具体输入的行检查器。

\subsection{从单行契约到矩阵与模式}
先看投影方向。若 $\rho$ 满足式 \eqref{eq:exact} 且 $x=E_R(H)$，
由于所有路线流非负，任何非树边在总流中为零时，每条路线在该边上的流也为零。
因此各条路线均有树支持。现有 \code{routeRow_root}、\code{routeRow_ancestor}
与 \code{routeRow_counts} 给出根出发一次、祖先闭合及访问数／成本的对应。
父边精确流等式给出 $\sum_i\pi_R(\rho_i)(v)=H(v)$（$v\neq o$），
由非负性再得每行坐标界 $\pi_R(\rho_i)(v)\leq H(v)$。
所以规范化投影同时得到合法矩阵和逐行的 $\RV$；这一步也允许有符号费用与预算。

对满足矩阵条件的 $p_0,\ldots,p_{m-1}$，先逐行规范化，再用命题
\ref{prop:section} 选取 $r_i$。非根列和不变，且
\[
 \sum_iF_{r_i}=\sum_iE_R(p_i)=E_R(H).
\]
由式 \eqref{eq:degrees} 和 \code{exactFlow_implies_coverage} 得到精确需求覆盖；
每行成本等式给出逐路线预算。选取顺序不影响上述等式。
这也是现有 \code{matrix_implies_exact} 的组合方式，新增数学规格明确的是其单行往返接口。

若把规范化矩阵行按完整向量相等分组，
$z(p)=\#\{i:p_i=p\}$，则
\[
 \sum_pz(p)=m,\qquad \sum_pz(p)p(v)=H(v)\quad(v\neq o).
\]
反向可把 $p$ 重复 $z(p)$ 次，从而恢复一个矩阵；
这只确定一个行多重集，不恢复原路线索引。
\code{family_iff_counts} 和 \code{matrix_iff_pattern_counts}
已证明相应存在性结论。
其中使用有限类型基数等价恢复 $\operatorname{Fin}(m)$ 的证明，不能直接当成保留原索引的执行算法。

\section{两种严格的信息损失}

\begin{example}[不同访问顺序投影到同一行]
取树边 $\{o,h\},\{h,a\},\{h,b\}$，完整路线为
\[
 r=[o,h,a,h,b,h,o],\qquad r'=[o,h,b,h,a,h,o].
\]
二者都是合法树支持路线，且 $r\neq r'$，但在坐标顺序 $(o,h,a,b)$ 下
\[
 \pi_R(r)=\pi_R(r')=(0,1,1,1),\qquad
 N_r=N_{r'}=(0,3,1,1),\qquad
 C_s(r)=C_s(r')=3s(h)+s(a)+s(b).
\]
因此投影不是单射。允许忽略访问顺序，恰恰因为当前约束只读取流、访问次数与可加费用；
若增加时间窗、访问先后或随历史变化的费用，就必须另行验证是否仍可丢弃顺序。
\end{example}

\begin{example}[总流不能检查一个给定分解的逐路线预算]\label{ex:aggregate}
取树边 $\{o,h\}$ 和 $\{h,a\},\{h,b\},\{h,c\},\{h,e\}$。
令 $s(h)=s(o)=0$，$s(a)=s(b)=2$，$s(c)=s(e)=1$，$m=2,T=3$。
每条路线依次访问分给它的两个叶子，每次在叶子之间经过 $h$。
在坐标顺序 $(o,h,a,b,c,e)$ 下，两种逐路线矩阵为
\[
 P_{\mathrm{good}}=
 \begin{pmatrix}0&1&1&0&1&0\\0&1&0&1&0&1\end{pmatrix},\qquad
 P_{\mathrm{bad}}=
 \begin{pmatrix}0&1&1&1&0&0\\0&1&0&0&1&1\end{pmatrix}.
\]
每行都满足根出发一次与祖先闭合；两矩阵列和均为
$H=(0,2,1,1,1,1)$，故总流均为 $E_R(H)$。
两种分解都使每个叶子访问一次，$h$ 总共访问 $2+4=6$ 次；
不是把 $h$ 的需求误写成 $2$ 或 $4$。
两种分解总成本均为 $6$，但逐路线成本分别为
\[
 (3,3)\quad\text{与}\quad(4,2).
\]
只有前者满足逐路线预算。
\end{example}

令 $\mathcal D$ 为满足结构、精确列和的这类行矩阵集合，
$A(P)=\sum_i p_i$ 为加总映射，$B(P)$ 表示所有行都满足预算。
例 \ref{ex:aggregate} 严格证明：不存在只依赖 $A(P)$ 的谓词 $b$，
使 $B(P)\leftrightarrow b(A(P))$ 对这里所有给定分解 $P\in\mathcal D$ 成立。
这是\emph{见证有效性}的信息损失。

它不证明固定输入 $(H,s,T,m)$ 的\emph{存在性可行性}不能判定。
该输入事实上可行，因为 $P_{\mathrm{good}}$ 已给出见证；
现有完整搜索就是从同一输入寻找合法分组。
因此不要把“有好分解也有坏分解”误写为“同一实例既可行又不可行”。

更强的既有分离结果位于
\src{FixedFlowObstruction}{\code{packing_feasible_but_fixed_flow_infeasible}}
和 \src{CopyLiftObstruction}{\code{physical_cuts_pass_but_copy_cuts_obstruct}}：
前者区分可行访问装箱与指定物理流的可实现性，后者区分物理聚合切割与副本容量切割。
它们使用各自的具体模型，不能直接当作上面这棵树的另一种分解结论。
本文仅引用既有源码中的结果，不将其历史验证状态重新表述为本次新运行。

\section{检查器契约与实现边界}

\subsection{模式压缩只有已声明的接受保持}
已有稳定压缩按\emph{完整行向量相等}合并重数，删除零重数行，
保留首次正重数出现的行顺序。\src{TreeBudgetCertificateCompression}{
\code{compress_checked}} 证明
\[
 \CC(C)=\mathrm{true}\quad\Longrightarrow\quad
 \CC(\operatorname{compress}(C))=\mathrm{true}.
\]
同时已有重数、加权列和、无重复性、幂等性和长度不增结论。
不能对任意 $C$ 改写为布尔值恒等式：原检查器会检查一个重数为零的不合法行，
压缩后该行消失，可能从拒绝变为接受。
源码回归
\src{TreeBudgetCertificateCompressionExamples}{
\code{checker_equality_is_not_unconditional}}
已经明确展示此边界。

一个合法矩阵经规范化、转成全为重数 $1$ 的证书，再压缩，必定仍被接受；
但反向展开只保证可获得合法矩阵，不保证回到原始条目表。
完整向量相等还包括根坐标；任意输入若根坐标不同，现有压缩不会擅自先归零后合并。
合法行本来已满足根坐标为零。

\subsection{哪些对象可以执行}
\small
\begin{longtable}{@{}p{.40\textwidth}p{.54\textwidth}@{}}
\toprule 组件 & 当前接口与本稿的边界\\\midrule
\endfirsthead
\toprule 组件 & 当前接口与本稿的边界\\\midrule
\endhead
\code{routeRow}、\code{normalize}、可执行 \code{cost} &
有限路线计数、根归零与有限和有明确计算定义；不因此得到编译程序的时间界。\\[4pt]
\code{single_row_realization}、\code{matrix_implies_exact} &
已存在的见证存在性证明；矩阵重建使用选择。
本文的 \code{realize} 是数学选取，不是已经提取的算法。\\[4pt]
\code{prevOfNext}、\code{toPairing}、\code{pairingOfRoutes} &
现有恢复定义使用 \code{noncomputable}／经典选择。
任务前驱唯一不等于源码已提供可执行前驱恢复程序。\\[4pt]
有限 next \code{checkCandidate} &
已有可执行检查器，通过直接副本流计算绕过经典恢复；检查给定候选，不搜索所有候选。
保留有限有效性定义的预算守卫与网格限制。\\[4pt]
模式 \code{checkCertificate}、\code{compressCertificate} &
已有有限检查和稳定压缩；接受保持及大小界已有证明。
函数式类型与有限求和的可执行性不等于已完成所有向量／位表示的成本分析。\\[4pt]
\#22 候选列表对应 &
属于另一实现接口：要求元素、顺序、出现次数的列表等号。
本稿的存在性和见证合法性桥接不要求所有未来算法返回同一个首见证。\\
\bottomrule
\end{longtable}
\normalsize

\section{后续 Lean 第一包的具体声明}

建议把下列三个小目标作为下一次形式化的边界，先不构造覆盖所有表示的大框架。
以下为\emph{规格}，不是本次新增定理清单。
\begin{enumerate}
 \item \textbf{共同范围桥接。}
 在命题 \ref{prop:bridge} 的完整假设下，组合 next、精确路线、矩阵与模式检查器的存在性等价。
 根出度等式和无自环应由树／度数条件推导，不作为重复且可能矛盾的输入。
 \item \textbf{见证级单行截面。}
 对满足 $\RV$ 的规范化行，证明存在 $r$ 使式 \eqref{eq:section} 成立。
 优先以“存在 $r$ 且投影等于 $p$”表达；若引入选择函数，名称和文档需注明非计算性。
 另给不要求规范化的版本 $\pi_R(r)=\norm(p)$，固定根坐标的边界。
 \item \textbf{严格信息损失。}
 形式化两条不同路线具有同一投影的小例；将其与给定分解的预算不能由总流检查的小例并列。
 保留已有压缩反例，避免把接受保持升级为任意输入的检查结果相等。
\end{enumerate}
单行重建不依赖服务费非负，共同 next 桥接仍按现有公开定理采用 $s\geq0,T>0$。
全矩阵与压缩表的反向展开可在后续见证级接口中规定具体索引与枚举顺序；
本包不提前声称编号、路线访问顺序、副本标号或缓存顺序恢复。

\section*{源码依据与状态说明}
全部定理名在本文读取的本地 UAV 源码中核对，引用链接固定到
\href{https://github.com/huam54925-cpu/uav-routing-certification/commit/5114dc8efc5b5082aefd4640b403285d129db2b2}{
\code{5114dc8}} 的研究代码基线。
\href{https://github.com/huam54925-cpu/uav-routing-certification/commit/06077788fc4a211126de98db7c75b88c9ef20db7}{
\code{0607778}} 是之后的状态文档更新；本稿没有把它当作新的证明构建。
相关议题为
\href{https://github.com/huam54925-cpu/uav-routing-certification/issues/21}{Issue \#21}
与 \href{https://github.com/huam54925-cpu/uav-routing-certification/issues/22}{Issue \#22}。

本稿只提供可独立放入 Overleaf 的 \code{.tex} 源码，建议选择 XeLaTeX。
未在本地编译 LaTeX，未生成 PDF；纸面命题的后续 Lean 编译、公理审计和回归结果须另行记录。
本稿没有进行文献新颖性判断，也没有推出强 NP 困难性、最优压缩或完整运行时间界。

\end{document}
